% Source: https://artofproblemsolving.com/community/c3h1222130_easy_geometry_help
% Stable identity: aops-topic-1222130-q01
Yes. For a convex $n$-gon and any interior point $X$, the segments to the vertices cut it into $n$ triangles, with total angle $180n^\circ$. The angles at $X$ fill a full turn, $360^\circ$, and the remaining angles of the triangles make up exactly the interior angles of the polygon (each interior angle is split by one segment into two triangle angles).

\begin{asy}
usepackage("amsmath");
size(3.5cm);
defaultpen(fontsize(10pt));
pair[] P={(0,0),(4,-0.5),(5.5,2),(3.5,4),(0.5,3.2)}; pair X=(2.5,1.6);
draw(P[0]--P[1]--P[2]--P[3]--P[4]--cycle,linewidth(1)); for(int i=0;i<5;++i) draw(X--P[i],blue); dot(X); label("$X$",X,S);
label("\mbox{}",(min(currentpicture,user=true)+max(currentpicture,user=true))/2,black); // keeps the text after the figure black when LaTeX typesets the labels inline
\end{asy}

So the angle sum is $180n-360=180(n-2)$ degrees, the usual formula. One need not draw triangles from a vertex; any point from which every vertex is visible works (for non-convex polygons a point must be chosen with this property, which may not exist).
